\newpage
\chapter{The Binomial Coefficient}

In this short section, we will cover the definition and properties of the \emph{binomial coefficient}.

\[
    \binom{n}{k} := \frac{n!}{k!(n-k)!},
\]

or as a product:

\[
    \prod_{j = 1}^{n} \frac{n-k+j}{j},
\]

with \(n\) and \(k\) as natural numbers (for the moment).

\subsubsection*{Properties}

\begin{itemize}
    
    \item \( \binom{n}{k} =  \binom{n}{n - k} \)
    \item \( \binom{n}{k - 1} = \binom{n}{k} \frac{n - k}{k + 1} \)
    \item \( \sum_{k = 0}^{n} \binom{n}{k} = 2^n \)
    \item \( \sum_{k = 0}^{n} \binom{m}{k} = \binom{n + 1}{m + 1} \)
    \item \( \binom{n}{0} = 1 = \binom{n}{n} \)
    \item \( \binom{n}{k} + \binom{n}{k + 1} = \binom{n + 1}{k + 1} \)
    \item \( \binom{n}{k} = {(-1)}^k \binom{k - n - 1}{k} \)

\end{itemize}

\subsection{The Binomial Theorem}

\[
    (a + b)^n =\sum_{k = 0}^{n} \binom{n}{k}a^{n - k}b^k
\]

\textbf{Proof:}

We will prove that the equality holds for any \(n\).

For \(n\), we have \( (a + b)^n = \sum_{k = 0}^{n} \binom{n}{k} a^{n - k} b^k \).

For \(n + 1\), we get:

\[
    (a + b)^{n + 1} =\sum_{k = 0}^{n + 1} \binom{n + 1}{k}a^{n + 1 - k}b^k.
\]

Now, let us manipulate the expression to demonstrate our theorem.

Assume that the statement is true for some integer \(n \ge 0\).

\[
    (a+b)^n = \sum_{k=0}^n \binom{n}{k} a^{n-k} b^{k}.
\]

Now, for some greater \(n\),

\[
    (a+b)^{n+1} = (a+b)(a+b)^n.
\]

Using the induction hypothesis:

\[
    (a+b)^{n+1} = (a+b) \sum_{k=0}^n \binom{n}{k} a^{n-k} b^{k}.
\]

We distribute:

\[
    (a+b)^{n+1} = \sum_{k=0}^n \binom{n}{k} a^{\,n+1-k} b^{\,k}
    + \sum_{k=0}^n \binom{n}{k} a^{\,n-k} b^{\,k+1}.
\]

In the second sum, change the index of summation by setting \(j = k + 1\), then rename \(j\) back to \(k\):

\[
    \sum_{k=0}^n \binom{n}{k} a^{\,n-k} b^{\,k+1}
    = \sum_{k=1}^{n+1} \binom{n}{k-1} a^{\,n+1-k} b^{\,k}.
\]

Now, combine both sums:

\[
    (a+b)^{n+1} = \sum_{k=0}^{n+1} 
    \left( \binom{n}{k} + \binom{n}{k-1} \right)
    a^{n+1-k} b^{k},
\]

where we use the convention \(\binom{n}{k} = 0\) for \(k < 0\) or \(k > n\).

By \emph{Pascal’s Identity}:

\[
    \binom{n}{k} + \binom{n}{k-1} = \binom{n+1}{k}.
\]

Therefore,

\[
    (a+b)^{n+1} = \sum_{k=0}^{n+1} \binom{n+1}{k} a^{\,n+1-k} b^{\,k}.
\]

\QED

\subsection{Pascal's Triangle}

This is a triangle built by the formula \(\binom{n + 1}{k + 1} = \binom{n}{k} + \binom{n}{k + 1}\).

\begin{center}
    \begin{tabular}{>{$n=}l<{$\hspace{12pt}}*{13}{c}}
    0 &&&&&&&1&&&&&&\\
    1 &&&&&&1&&1&&&&&\\
    2 &&&&&1&&2&&1&&&&\\
    3 &&&&1&&3&&3&&1&&&\\
    4 &&&1&&4&&6&&4&&1&&\\
    5 &&1&&5&&10&&10&&5&&1&\\
    6 &1&&6&&15&&20&&15&&6&&1
    \end{tabular}
\end{center}
\vspace{1cm}

Here is the triangle, but with the binomial coefficients:

\begin{center}
    \begin{tabular}{>{$n=}l<{$\hspace{12pt}}*{13}{c}}
    0 &&&&&&&\(\binom{0}{0}\)&&&&&&\\
    1 &&&&&&\(\binom{1}{0}\)&&\(\binom{1}{1}\)&&&&&\\
    2 &&&&&\(\binom{2}{0}\)&&\(\binom{2}{1}\)&&\(\binom{2}{2}\)&&&&\\
    3 &&&&\(\binom{3}{0}\)&&\(\binom{3}{1}\)&&\(\binom{3}{2}\)&&\(\binom{3}{3}\)&&&\\
    \end{tabular}
\end{center}

It can be used with the Binomial Theorem to get the coefficients of a polynomial of degree \(n\).

\textbf{Example:}

\[
    {(a + b)}^2 = 1 a^2 b^0 + 2 a^1 b^1 + 1 a^0 b^2,
\]

or

\[
    {(a - b)}^2 = 1 a^2 b^0 - 2 a^1 b^1 + 1 a^0 b^2.
\]
